% ============================================================================
% SECTION 6: CONCLUSION
% ============================================================================
\section{Conclusion}
\label{sec:conclusion}

This paper presented a behavioral characterization of how CNN and Vision Transformer architectures differ under adversarial attack. Moving beyond simple robustness comparisons, we examined these differences through transfer attack analysis, gradient characterization, and layer-wise feature degradation studies.

Our analysis establishes a statistically significant but modest \emph{absolute} robustness advantage for the CNN under full-test AutoAttack evaluation (35.74\% vs 32.94\%, McNemar $p < 10^{-11}$), together with a conditional decomposition that qualifies it: the PGD-10 gap is fully accounted for by the clean-accuracy gap (conditioned fooling 52.15\% vs 52.33\%), the per-model AutoAttack conditional ordering slightly favors the ViT, and yet on the jointly solved sample subset the CNN's advantage persists under both attacks---so the headline comparison depends on the view, and we report all three. The WideResNet baseline's much higher robustness further shows that capacity, training data, and training recipe jointly matter beyond architectural choice. A key methodological finding concerns transfer measurement: the apparent cross-architecture asymmetry depends on the conditioning protocol. Raw rates manufacture an 8.3-point asymmetry out of the targets' unequal clean accuracy; the target-correct conditioned metric shows statistical parity; and the stricter both-correct and successful-source protocols reveal a significant 5.3-point CNN$\rightarrow$ViT advantage on jointly solved samples. Transfer studies should therefore state and justify their conditioning protocol explicitly---none of these views licenses interpreting raw-rate asymmetries architecturally. Gradient analysis with scale-invariant measures shows CNN sensitivity is modestly more concentrated---a difference that survives a native-resolution ViT control and per-sample paired tests---while cross-sample gradient alignment is low for both architectures, moderately higher for the ViT in absolute cosine, and near zero in signed cosine for both. We also identify a characteristic degradation profile whereby adversarial perturbations cause directional deviation in ViT feature representations that accumulates by mid-network (cosine similarity falling from 0.994 to 0.955) and then plateaus, despite feature-norm changes below 1.5\%---in contrast to the CNN's monotonic depth profile.

This work characterizes gradient structure alongside cross-architecture transfer behavior, dissociating which gradient properties are architectural (cross-sample alignment) and which are shaped by adversarial training (magnitude and sparsity ordering), and offering methodological guidance in addition to performance comparison. Our systematic layer-wise analysis characterizes how adversarial perturbations propagate through ViT layers, and we distill practical guidance for evaluating and deploying mixed-architecture ensembles. We will release our analysis pipeline to enable future work extending these findings to other architectures and datasets.

% ----------------------------------------------------------------------------
% Closing
% ----------------------------------------------------------------------------
Understanding \textit{how} architectures behave differently under adversarial attack---not just \textit{how much}---is essential for developing targeted defenses. Our analysis shows that the behavioral differences between CNNs and ViTs correlate with scale-invariant differences in gradient structure, particularly how sensitivity is distributed across input components---an effect we verified is not an artifact of gradient scale or input preprocessing---though disentangling architectural effects from capacity and training-recipe factors requires the capacity-controlled follow-ups we outline. Cross-sample gradient alignment is low for both adversarially trained models yet remains moderately higher for the ViT (0.052 vs 0.038, 1.36$\times$); measuring the clean-trained counterparts shows nearly identical levels (0.050 vs 0.030), identifying the alignment gap as an architectural property largely independent of the training regime---while gradient sparsity, whose architecture ordering adversarial training reverses, emerges as the training-shaped characteristic. These dissociations open new questions about the interaction between defense methods and architectural inductive biases. We hope this behavioral perspective guides future research toward more principled defense design.
